% Generated draft: requires mathematical and source review.
\documentclass[11pt]{article}
\usepackage[margin=25mm]{geometry}
\usepackage{fontspec}
\setmainfont{Latin Modern Roman}
\usepackage{amsmath,amssymb,mathtools,unicode-math}
\setmathfont{Latin Modern Math}
\usepackage{enumitem,needspace,xurl,hyperref}
\setlength{\parindent}{0pt}
\setlength{\parskip}{5pt}
\setlength{\emergencystretch}{3em}
\newcommand{\N}{\mathbb{N}}
\newcommand{\Z}{\mathbb{Z}}
\newcommand{\Q}{\mathbb{Q}}
\newcommand{\R}{\mathbb{R}}
\newcommand{\C}{\mathbb{C}}
\newcommand{\F}{\mathbb{F}}
\newcommand{\A}{\mathbb{A}}
\newcommand{\PP}{\mathbb{P}}
\newcommand{\E}{\mathbb{E}}
\newcommand{\eps}{\varepsilon}
\newcommand{\dd}{\,\mathrm{d}}
\begin{document}
\section*{UnsolvedMath ID: 30002296}
\textbf{Source problem code:} catalogue entry
\subsection{Definitions and mathematical statement}
A family \((f_i)_{i\in I}\) in a Hilbert space \(H\), indexed by a countable set \(I\), is a tight frame if there exists a constant \(A>0\) such that
\[
\sum_{i\in I}|\langle f,f_i\rangle|^2=A\|f\|_H^2
\qquad\text{for every }f\in H.
\]
When \(A=1\), it is a Parseval frame. A wavelet generator is a function whose prescribed dilations and translations contribute to the associated wavelet system; the number of generators counts these functions, rather than all members of the generated system.

The supplied problem asks: \emph{Obtain estimates for the number of generators needed by multivariate tight wavelet frames.}

To describe the counting target without choosing missing wavelet hypotheses, let \(D\) denote specified wavelet data: an ambient Hilbert space \(H_D\), the dilation and translation rules, the admissibility requirements on generators, and any prescribed refinable function or mask. Let \(\mathcal A_D(r)\subseteq H_D^r\) be the admissible ordered families of \(r\) generators, and let \(\mathcal W_D(\Psi)\) be the wavelet system generated by \(\Psi\). One possible precise counting quantity, once these data and the intended optimization have been fixed, is
\[
r_{\min}(D)=\inf\bigl\{r\in\mathbb Z_{\geq1}:\exists\Psi\in\mathcal A_D(r)\text{ such that }\mathcal W_D(\Psi)\text{ is a tight frame for }H_D\bigr\},
\]
with \(\inf\varnothing=+\infty\). Bounds on this quantity would express a minimum-generator interpretation of the request. This is a conditional formulation scheme, not an assertion that the source specifies this particular optimization.

The supplied statement does not specify the dimension range, dilation, translation lattice, ambient function space, generator regularity or support, prescribed refinement data, or parameters on which estimates should depend. It also does not determine whether the requested estimates concern arbitrary tight wavelet frames or constructions subject to the unitary extension principle (UEP). Consequently, a unique quantified bound problem cannot be recovered from the supplied statement alone.

The supplied dated curation labels the problem partially solved and reports that Charina--Putinar--Scheiderer--Stockler resolve the bivariate UEP existence question, while later work improves generator-count upper bounds and constructions without providing a complete multivariate estimate. The hypotheses and scope of those results are not supplied here; this assessment is not an independently verified mathematical conclusion.
\subsection{Short English statement}
Estimate how many generating functions are needed to construct multivariate tight wavelet frames. The source leaves the admissible systems, construction constraints, and parameters of the desired bounds unspecified, so a precise bound requires further source clarification. The supplied partially solved label is dated curation.
\subsection{Sources}
\begin{itemize}
\item Local UnsolvedMath registry, permanent ID \(30002296\), source problem code \texttt{catalogue entry}, title: \emph{Generator Bounds for Multivariate Tight Wavelet Frames}. The supplied original, clean, and current statements agree. The record gives proposed year 2013 and identifies no individual proposer.
\item \emph{Structured Function Systems and Applications} (2013), Oberwolfach Reports, pp.\ 579--655, DOI 10.4171/owr/2013/11. Supplied source report and literature-source label: ``Oberwolfach Report 2013.'' \url{https://doi.org/10.4171/owr/2013/11}.
\item Charina--Putinar--Scheiderer--Stockler, \emph{Multivariate tight wavelet frames}. Supplied literature reference: \url{https://doi.org/10.1007/s00365-013-9191-5}.
\item Charina--Putinar--Scheiderer--Stockler, \emph{Part II}, as labeled in the supplied record. \url{https://doi.org/10.1016/j.acha.2014.09.003}.
\item Literature assessment checked at \texttt{2026-08-21}: the supplied curation labels the problem partially solved, attributes resolution of the bivariate UEP existence question to Charina--Putinar--Scheiderer--Stockler, and reports subsequent improvements in generator-count upper bounds and constructions without a complete multivariate estimate. This is supplied dated curation derived from the Oberwolfach material, not independently verified here; time-sensitive status and cited claims require independent verification.
\item The registry marks the statement ``exact'' and reports that its mathematical content was checked against the official Oberwolfach report, with notation, context, and question form normalized without changing the source-level claim. This verification description is supplied provenance, not a verification performed for this draft.
\end{itemize}
\end{document}
